\section{Discussion}
\label{sec:nmi-discussion}
% Alias label so the reused appendix (SI) cross-references resolve in this variant.
\label{sec:discussion}

\paragraph{Intelligence is not the bottleneck; reliability is.}
The load-bearing result is that Direct (GPT-5.4), a one-paragraph prompt, already separates
rejected from accepted submissions on its own (AUROC \aurocNaiveFull), indistinguishably from
AIPR ($\Delta$AUROC \aurocDiffFullNaive, \aurocDiffPrel). Raw judgment is not what this
deployment is shortest of; reliability is---the bare prompt's score swinging \naiveRunSD{}
points between identical runs against AIPR's \fullRunSD{}. The open problem is not building a
model that \emph{can} judge but turning that capability into an instrument a reviewer can rely
on: one that returns the same verdict across runs, grounds its claims in evidence, and surfaces
an anchored review instead of a number. That is an engineering problem, and the one AIPR
addresses; the field's effort is better spent on the processing layer---reliability, grounding,
interaction---than on a more intelligent base model the task does not appear to need.

\paragraph{Compressed scores are a feature, not a defect.}
AIPR's production scores cluster high, near the venue bar, and move little between runs. Such
clustering is a known behaviour of scalar LLM judges \citep{girrbach2025referencefree}, and for
a deployed grader reliability rather than spread is the property that governs whether the score
can be trusted \citep{gu2024surveyllmjudge,schroeder2024trustllm}. We trade spread for stability
deliberately: a stable estimate near a reference value is easier to act on than a wider-ranging
but noisier one, and a low point estimate carried with a wide interval fails as a deployment
object because a reader anchors on the point and under-weights the uncertainty around it
\citep{hofman2020inferential}. We do not claim compression improves discrimination (AIPR's AUROC
is not higher than the bare prompt's) or reviewers' actual decisions, which we have not measured;
the claim is about the estimator and how it is read.

\paragraph{Policy, contamination, and an obvious next step.}
This reframes the policy question. Recent randomized evidence shows AI feedback can improve
review \emph{text} \citep{thakkar2026feedback}; our complementary result is that a first-pass
\emph{score} can be criterion-valid against human outcomes. Neither licenses uncontrolled use:
undisclosed AI-assisted review \citep{russolatona2024lottery} and adversarial hidden prompts in
manuscripts \citep{lin2025hidden} are exactly why assistance must be a controlled, disclosed,
auditable system rather than reviewers pasting confidential manuscripts into general tools. The
central threat to our own claim is contamination---that the grading model memorized the
outcome---which we control by cohort choice: GPT-5.4's August~2025 cutoff \citep{openai2026gpt54}
precedes the January~2026 release of the decisions, bounded further by an arXiv-before-cutoff
sensitivity split (SI). The deployed score's lightly weighted citation dimension was
uninformative this run for a technical reason in its audit; adding a grounded, search-backed
citation check is the natural next step, and matters given fabricated references in accepted
papers \citep{neurips2025citations}.

\paragraph{Limitations.}
\label{sec:limits}%
Several limits bound the claim, and we state them plainly. The reject class is weak
\emph{relative to the venue bar}, not absolutely, and the human decision is itself noisy
\citep{cortes2021inconsistency,beygelzimer2023arbitrary}; we mitigate by also validating against
the mean rating, but no analysis exceeds its ground truth, and that truth may itself be partly
LLM-assisted \citep{liang2024monitoring}. The scale rests on the cheap-model bridge, weakest at
the low end, which is why we report the deployable flag on the production score directly. The
evidence is one venue in one field; a within-venue temporal cohort and a second-field
replication are planned follow-ups. And the scorer is proprietary: the data, labels, and
analysis are released and fully reproducible, but the manuscript-to-score function is audited
through its outputs rather than re-implemented, with prompts and configuration available
confidentially to editors. We do not claim the score predicts acceptance, ranks strong papers,
or should drive any decision without a human in the loop. The validated action is prioritizing
manuscripts for human attention, not automated rejection.
